\documentclass[12pt,a4paper]{article}
\usepackage[utf8]{inputenc}
\usepackage[spanish]{babel}
\usepackage{geometry}
\geometry{top=2.5cm, bottom=2.5cm, left=3cm, right=3cm}
\usepackage{graphicx}
\usepackage{hyperref}
\usepackage{xcolor}
\usepackage{fancyhdr}
\usepackage{amsmath}

\definecolor{unabblue}{HTML}{003865}
\definecolor{unaborange}{HTML}{E85C0B}
\definecolor{codegray}{rgb}{0.95,0.95,0.95}
\definecolor{codeorange}{rgb}{0.8,0.3,0.0}

\pagestyle{fancy}
\fancyhf{}
\fancyhead[L]{\textcolor{unabblue}{\textbf{Ciencia de Datos - UNAB}}}
\fancyhead[R]{\textcolor{unaborange}{\textbf{Taller Semana 4}}}
\fancyfoot[C]{\thepage}

\begin{document}

\begin{center}
    {\huge \textcolor{unabblue}{\textbf{Taller 4: Data Janitor 3000 (Limpieza y SQL)}}} \\[0.5cm]
    {\large Docente: Dudbil Olvasada Pabon Riaño} \\[0.2cm]
    \rule{\textwidth}{1pt}
\end{center}

\vspace{0.5cm}

\section*{\textcolor{unabblue}{1. Objetivo del Taller}}
Enfrentar la realidad de los datos del mundo real: la suciedad y la inconsistencia. El estudiante aplicará técnicas de imputación estadística usando Pandas y ejecutará cruces lógicos simulando operaciones de bases de datos relacionales (SQL Joins).

\section*{\textcolor{unabblue}{2. Instrucciones Generales}}
Desarrolle los retos en un Jupyter Notebook. Para este taller necesitará generar los \texttt{DataFrames} de juguete propuestos usando diccionarios. Suba el documento a su repositorio documentando claramente por qué tomó cada decisión de limpieza.

\vspace{0.5cm}
\noindent\fbox{%
    \parbox{\dimexpr\textwidth-2\fboxsep-2\fboxrule\relax}{%
        \textbf{\textcolor{unabblue}{Misión Analítica: El Rescate del CRM}}\\[0.2cm]
        Usted es el ``Data Janitor" de una aerolínea. El equipo de reservas migró la base de datos antigua a un nuevo sistema, pero en el proceso, se generaron valores nulos (\texttt{NaN}), edades ilógicas y clientes duplicados. Si los modelos de Machine Learning entran en producción con esta basura, la empresa perderá millones.
    }%
}
\vspace{0.5cm}

\subsection*{Reto 1: Detectando la Basura Oculta (20\%)}
Cree el siguiente DataFrame en Pandas usando este código:
\begin{verbatim}
import pandas as pd
import numpy as np

data = {
    'cliente_id': [101, 102, 103, 104, 105, 106, 101],
    'edad': [25, 30, np.nan, 800, 45, 22, 25],
    'genero': ['M', 'F', 'F', 'M', 'O', np.nan, 'M'],
    'millas': [1500, 2000, 50, 10000, np.nan, 300, 1500]
}
df = pd.DataFrame(data)
\end{verbatim}

\begin{enumerate}
    \item Utilice los métodos \texttt{.isnull().sum()} o \texttt{.info()} para contar exactamente cuántos valores nulos (\texttt{NaN}) tiene cada columna. Imprima el reporte.
    \item \textbf{Análisis Crítico:} Observe el valor de la edad del cliente 104. ¿Es posible tener 800 años? Describa en una celda de texto cómo esta ``anomalía semántica" (que no es un \texttt{NaN}) afectaría el promedio (media) de edades si la pasamos por alto.
\end{enumerate}

\subsection*{Reto 2: Cirugía de Datos: Imputación (30\%)}
No podemos simplemente eliminar filas porque perderíamos información valiosa de los clientes. Debemos ``Imputar" (rellenar inteligentemente).
\begin{enumerate}
    \item \textbf{Imputar Variables Numéricas:} Rellene el valor nulo de la columna \texttt{millas} utilizando la \textbf{mediana} de esa misma columna. (Pista: investigue el método \texttt{.fillna()}).
    \item \textbf{Imputar Variables Categóricas:} Rellene el valor nulo de la columna \texttt{genero} utilizando la \textbf{moda} (el valor más frecuente) de la columna.
    \item \textbf{Tratar Valores Absurdos:} Reemplace el valor absurdo de \texttt{800} en la edad por un \texttt{NaN} usando el método \texttt{.replace()}, y posteriormente impútelo usando la media de la edad. Muestre el DataFrame resultante.
\end{enumerate}

\subsection*{Reto 3: Desduplicación (10\%)}
Durante la migración, algunos registros de clientes se insertaron dos veces. 
\begin{enumerate}
    \item Utilice el método \texttt{.drop\_duplicates()} de Pandas para eliminar filas idénticas.
    \item ¿Cuántas filas quedaron en total después de eliminar los duplicados?
\end{enumerate}

\subsection*{Reto 4: SQL en Pandas (El Arte del Join) (40\%)}
Los datos demográficos ya están limpios, pero el área de fidelización tiene otra tabla separada con el nivel (Tier) de cada cliente.
Cree este segundo DataFrame:
\begin{verbatim}
tier_data = {
    'cliente_id': [101, 102, 104, 105, 999],
    'nivel': ['Gold', 'Silver', 'Platinum', 'Silver', 'Gold']
}
df_tier = pd.DataFrame(tier_data)
\end{verbatim}

\begin{enumerate}
    \item Realice un \textbf{INNER JOIN} (cruce interno) entre su DataFrame limpio (\texttt{df}) y el nuevo DataFrame (\texttt{df\_tier}) utilizando la columna \texttt{cliente\_id} como llave primaria (\textit{Key}). Utilice el método \texttt{pd.merge()}.
    \item Realice un \textbf{LEFT JOIN} (cruce por la izquierda). Mantenga todos los clientes de su \texttt{df} original, cruzando con \texttt{df\_tier}. 
    \item \textbf{Análisis Crítico:} Observe el resultado del Left Join. ¿Qué sucede con los clientes 103 y 106 en la columna \texttt{nivel}? ¿Por qué aparece un \texttt{NaN} ahí, y por qué el cliente 999 de la tabla de Tiers no aparece en absoluto? Explique la lógica de la teoría de conjuntos detrás de esto.
\end{enumerate}

\vspace{1cm}
\begin{center}
    \textit{\textcolor{gray}{`Tu modelo de Inteligencia Artificial solo es tan inteligente como los datos basura que le diste de comer (Garbage In, Garbage Out).'}}
\end{center}

\end{document}
